\section{Introduction}
\label{sec:intro}

Retrieval-augmented generation is only as good as the passages it retrieves,
and a chunk that has been cut out of its document frequently reads as
ambiguous on its own. Contextual retrieval attacks this at indexing time by
prepending to each chunk a short summary of the document it came from, which
measurably reduces retrieval failures \citep{anthropic2024contextual}. The
published recipe does not stop there: it also reranks the fused candidate list
with a cross-encoder.

Neither stage is cheap once the system leaves a datacenter. Enrichment costs a
language-model call per chunk, though this is paid once, at ingest, and can be
amortized over many queries. The reranker is paid on every query, one forward
pass per candidate, and on a laptop or a phone that is the cost that hurts:
sixty-four candidates through a 0.6B cross-encoder take tens of seconds, while
a dense-only first stage answers in tens of milliseconds. If the contextual
enrichment has already resolved the ambiguity the reranker was meant to fix,
the reranker can be dropped, and the whole pipeline becomes usable
interactively on a single device. This paper takes up that possibility and
tests it on a corpus of the kind an edge deployment actually meets.

Three questions organize the study. Is a reranker still needed once chunks are
contextualized? Can the enrichment itself be moved on-device, by distilling an
API teacher into a small local model, without giving up retrieval quality? And
which parts of the standard recipe, BM25 fusion, embedding dimension and the
language of the query instruction among them, actually help on a small,
homogeneous, non-English corpus? Published results report the recipe's two
components together, so the three questions cannot be answered from them
directly.

The corpus is 82 Russian business documents that came out of our own OCR
pipeline. It is deliberately small and stylistically uniform: a few hundred
documents that share vocabulary, layout and boilerplate, which is the setting an
edge deployment faces. Every one of the 3{,}134 chunks is labeled by an API
teacher, a 0.6B student is trained to reproduce the teacher's contexts, and
retrieval is evaluated with a purpose-built set of 234 queries carrying gold
chunks. All measurements are taken on a single Apple~M2 with MLX in fp16, the
same machine that would serve the queries.

Our contributions are the following.

\begin{enumerate}
  \item A contextual chunk dataset of 3{,}134 labeled chunks over 82 Russian
  documents, with category and table metadata, produced by an API
  teacher for USD~0.82 and released publicly (\S\ref{sec:data}).
  \item A 0.6B LoRA student that reproduces the teacher's contexts (word-level
  F1 improves from 0.189 to 0.537) and can regenerate the entire corpus
  on-device in 4.49~h at zero marginal cost (\S\ref{sec:distill}).
  \item An eight-arm retrieval evaluation, six configurations on 234 gold-chunk
  queries and all eight on a common stratified subset of sixty, showing that a
  reranker-free pipeline with distilled on-device context comes within 0.05
  recall@10 of the same index with a 0.6B cross-encoder at a factor of 1{,}322
  in median latency, and that reranking flat chunks is dominated by the
  contextual dense-only arm (\S\ref{sec:results-main}).
  \item A negative result for lexical fusion: on this homogeneous corpus BM25
  degrades every configuration, and neither the fusion weights, the rank
  constant nor the retrieval depth recover the loss (\S\ref{sec:ablations}).
  \item Two cheap levers for the on-device path, Matryoshka truncation to 512
  dimensions and a Russian query instruction, that together lift the student
  pipeline to the teacher baseline at half the index size
  (\S\ref{sec:ablations}).
  \item A stratified analysis of where distillation transfers and where it does
  not, and a break-even model that separates ingest cost from query cost
  (\S\ref{sec:strata}, \S\ref{sec:breakeven}).
\end{enumerate}
